\documentclass[a4paper,12pt]{article}

\usepackage[T2A]{fontenc}
\usepackage[utf8]{inputenc}
\usepackage[russian]{babel}

\usepackage{amsmath, amssymb}
\usepackage{geometry}
\usepackage{fancyhdr}
\usepackage{enumitem}

\geometry{margin=2cm}
\setlength{\parskip}{6pt}
\setlength{\parindent}{0pt}

\pagestyle{fancy}
\fancyhf{}
\rhead{7 класс}
\cfoot{\thepage}

\begin{document}

\begin{center}
    {\Large \textbf{Делимость. Простые числа.}}\\[0.5em]
\end{center}
    
\begin{enumerate}[label=\textbf{Задача \arabic*.}]
\item Какое наименьшее натуральное число не является делителем 50! ($n! = 1 \cdot 2\cdot3 \cdot\ldots\cdot n$)?
\item Докажите, что если $\ell$ - делитель числа n, то и число $\frac{n}{\ell}$ - делитель n.
\item Может ли n! оканчиваться пятью нулями?
\item Числа n и  $n^{2} + 2$  – простые. Докажите, что  $n^{3} + 2$  – также простое число.
\item В ряд выписаны числа от 1 до 25. Можно ли расставить между ними знаки «+» и «–» так, чтобы
значение полученного выражения было равно нулю?
\item Числа $1, 2, 3, …, 456$ выписаны в строчку по порядку. За ход разрешается менять местами любые два числа, стоящие через одно (например, 196 и 198). Можно ли с помощью таких ходов расположить все числа в обратном порядке? 
\item При каких натуральных n выражение $n^3 - n$ кратно 6?
\item Докажите, что для всякого натурального числа n $7^n +5$ кратно 6.
\item $x + 1$  делится на 3. Докажите, что  $1 + 7x$  делится на 3.
\item Докажите, что при любом натуральном $n > 5$ число  $n^2 - 8n + 15$  не делится на  $n - 4$ (воспользуйтесь тем, что $a^2 - 2ab + b^2 = (a-b)^2$). 
\item Найдите остаток числа $26^{36}$ при делении на 7.
\item Какая натуральная степень числа 2 оканчивается четырьмя одинаковыми цифрами?
\item Докажите, что число является квадратом целого числа тогда и только тогда, когда оно имеет нечетное число натуральных делителей.
\item Верно ли, что бесконечных чисел бесконечно много?

\end{enumerate}
\end{document}
